\subsection{Colored Noise (\texorpdfstring{$\mathbf n$}{n})}\label{subsec:interferencia}
The BS receives signals from the $K$ UEs it serves (target UEs), whose symbols form the vector $\mathbf s$. Signals received from non-target UEs interfere with data reception and are treated as noise. These signals, together with thermal noise, form the vector $\mathbf n$ in Equation~\ref{eq:systemmodel}. The interference makes the covariance matrix of $\mathbf n$ non-diagonal, and for this reason $\mathbf n$ is called colored noise \cite{zhao2025jsac}.
% FONTE: zhao2025jsac paper.md:33, 108-113, 577.

\subsubsection{Thermal Noise}\label{subsubsec:ruido-termico}
Thermal noise arises from the random motion of electrons in the conductors and devices of the receiver front end. This noise is modeled as additive white Gaussian noise (AWGN), which has a Gaussian distribution with zero mean and the same power spectral density at all frequencies \cite{haykin,tse-viswanath}. The noise at the $M$ antennas is described by an $M \times M$ covariance matrix, in which each element is the correlation between the noise of a pair of antennas. Thermal noise is independent at each antenna, so only the diagonal elements of this matrix are non-zero. Each diagonal element is the average thermal noise power at one antenna, $\sigma^2$, and the covariance matrix is $\sigma^2\mathbf I_M$ \cite{zhao2025jsac}.
% FONTE: haykin pp. 186, 225; tse-viswanath p. 29; zhao2025jsac paper.md:108-113.
